%% LyX 2.4.4 created this file.  For more info, see https://www.lyx.org/.
%% Do not edit unless you really know what you are doing.
\documentclass[english]{article}
\usepackage[T1]{fontenc}
\usepackage[latin9]{inputenc}
\usepackage{enumitem}
\usepackage{amsmath}
\usepackage{amsthm}
\usepackage{amssymb}
\usepackage{geometry}
\geometry{verbose,tmargin=1in,bmargin=1in,lmargin=1in,rmargin=1in}

\makeatletter
%%%%%%%%%%%%%%%%%%%%%%%%%%%%%% Textclass specific LaTeX commands.
\numberwithin{equation}{section}
\numberwithin{figure}{section}
\newlength{\lyxlabelwidth}      % auxiliary length 

%%%%%%%%%%%%%%%%%%%%%%%%%%%%%% User specified LaTeX commands.
\usepackage{fancyhdr}
\usepackage{lastpage}
\pagestyle{fancy}
\usepackage{enumitem}
\usepackage{ifthen}
%sets math fonts to be more readable
\everymath=\expandafter{\the\everymath\displaystyle}
%\DeclareMathSizes{10}{10}{10}{10}
\usepackage{multicol}
% Added by lyx2lyx
\setlength{\parskip}{\smallskipamount}
\setlength{\parindent}{0pt}

\makeatother

\usepackage{babel}
\begin{document}
\renewcommand{\author}{Dr.\,James Ashe}

\newcommand{\classNum}{331}

\newcommand{\classSec}{A}

\newcommand{\nameType}{Name:}

\newcommand{\examType}{Exam 4 Review}

\renewcommand{\date}{}

%%First page's header%%%%%%%%%%%%%%%%%%%%%%
\fancypagestyle{firststyle} %%header settings for first pagr
{
	\fancyhead[L]{MTH \classNum{}(\classSec{}) \author{}\\
	\textbf{\examType{}} \date{}}
	\fancyhead[C]{\textbf{\nameType}}
	\setlength{\headsep}{14pt}
}
%%%%%%%%%%%%%%%%%%%%%%%%%%%%%%%%%%%%%%%%%%

%%Next page's header%%%%%%%%%%%%%%%%%%%%%%%
\setlist{nolistsep}
\lhead{\classNum{}(\classSec{})} 
\chead{\textbf{\nameType}} 
\cfoot{\thepage{}~of~\pageref{LastPage}} 
\setlength{\headsep}{5pt} 
%%%%%%%%%%%%%%%%%%%%%%%%%%%%%%%%%%%%%%%%%%%

%%Various settings; see the preamble for other settings%%%%%
%%\setenumerate[0]{leftmargin=12pt,itemindent=0pt,labelwidth=0pt}
\renewcommand{\labelenumi}{\textbf{\arabic{enumi}.}}
\renewcommand{\labelenumii}{\textbf{\alph{enumii}.}}
\thispagestyle{firststyle}
%%%%%%%%%%%%%%%%%%%%%%%%%%%%%%%%%%%%%%%%%%

\section*{12.1: Vectors in the Plane}
\begin{itemize}
\item Basic vector operations and arithmetic: \textbf{17-20,} \textbf{39,
41, 43}

\begin{itemize}
\item Magnitude and unit vectors.
\item Adding, subtracting, scalars.
\item Parallelogram/Triangle rule.
\end{itemize}
\item Word problems: \textbf{45, 47}

\begin{itemize}
\item Describe various motions as vectors or as a combination of vectors.
\item Find a vector description given the direction and magnitude of a force. 
\item Find the magnitude and direction given a vector. 
\end{itemize}
\end{itemize}

\section*{12.2: Vectors in Three Dimensions}
\begin{itemize}
\item Graphing points in $\mathbb{R}^{3}$: \textbf{9-21}
\item Basic vector operations and arithmetic in $\mathbb{R}^{3}$ (similar
to 12.1 above): \textbf{35, 37}
\end{itemize}

\section*{12.3: Dot Products}
\begin{itemize}
\item Compute the dot product, $\mathbf{u\cdot v}$: \textbf{13-17}
\item Find the angle between two vectors (in 2 or 3 dimensions): \textbf{13-17}
\item Find $\mbox{proj}_{\mathbf{v}}\mathbf{u}$ and $\mbox{scal}_{\mathbf{v}}\mathbf{u}$:
\textbf{19-21, 25-27}
\item Word problems: \textbf{29-31}

\begin{itemize}
\item Find the \emph{work}, given vectors and/or magnitude and directions. 
\end{itemize}
\end{itemize}

\section*{12.4: Cross Products}
\begin{itemize}
\item Compute the cross product, $\mathbf{u\times v}$, and find an orthogonal
vector: \textbf{9-11, 25, 29}
\item Find torque: \textbf{33, 57}
\item Word problems: \textbf{37, 38, 57, 58}

\begin{itemize}
\item Describe a resulting force as the cross product of vectors.
\item Find the magnitude and direction of the force.
\end{itemize}
\end{itemize}

\section*{12.5: Lines and Curves in Space}

Continuous motions can be described using vector-valued functions
(also called position vector),\\
 $\mathbf{r}(t)=\left\langle x(t),y(t),z(t)\right\rangle $. This
is a parametric function in $\mathbb{R}^{3}$. 
\begin{itemize}
\item Concept questions about vector-valued functions: \textbf{1-8, 39-41}
\item Finding an equation of a line meeting given conditions: \textbf{9-13}
\end{itemize}

\section*{12.6: Calculus of Vector-Valued Functions}
\begin{itemize}
\item Find the derivative of a vector-valued function, $\mathbf{r}(t)$:
\textbf{7-15, 19, 31}

\begin{itemize}
\item Review the power, product, quotient, and chain rules. Also review
the derivative rules for the natural logarithmic, exponential, and
trigonometric functions. 
\end{itemize}
\item Find the definite and indefinite integral of $\mathbf{r}(t)$: \textbf{41-43,
49-53}

\begin{itemize}
\item Review the power and substitution rule. Also review the integration
rules for the natural logarithmic, exponential, and trigonometric
functions. You will not need to use integration by parts.
\end{itemize}
\item You will be provided all the necessary derivative and integration
rules (as they appear in the text).
\end{itemize}

\end{document}
